\section{What a wrong forward does to the smile}
\label{sec:fwdskew}

The chapter opened with the statement that a small forward error is a
large smile error.  Every tool needed to prove it is now on the
table, and the proof is a short computation with a memorable shape.

Set the scene precisely.  The market's quotes are what they are; the
error lives in the modeler's head.  The true forward is $F$, but the
smile is built at a misread forward $F'=Fe^{\delta}$ --- $\delta$ is
the relative error, a few basis points.  Every dollar quote is then
pushed through the Black formula at the wrong normalization: the
strike's log-moneyness is computed as $k'=\log(K/F')=k-\delta$, and
the price is divided by $DF'$ instead of $DF$.  The discount is read
correctly throughout --- the forward error is the subject.  The
question: how far does the re-implied volatility move, per unit of
$\delta$?

Take a call first.  Its dollar price, held fixed, is
$DF\,\Black(k,w)$ with the normalized Black call of Chapter~2
(eq.~\ref{eq:black}),
\[
\Black(k,w)=\PhiN(d_+)-e^k\PhiN(d_-),
\qquad
d_\pm=-\frac{k}{\sqrt w}\pm\frac{\sqrt w}{2},
\]
where $\PhiN$ and $\phin$ are, as throughout the book, the standard
normal distribution function and density.  The call's two partial
derivatives were computed once in Chapter~2's appendix, and this
chapter reuses them:
\[
\partial_k\Black=-e^k\,\PhiN(d_-),
\qquad
\partial_w\Black=\frac{\phin(d_+)}{2\sqrt w}.
\]
Differentiate the fixed dollar price along the misreading.  Three
things move at once: the prefactor $F$ (its log by $\dd\delta$), the
moneyness $k$ (by $\dd k=-\dd\delta$), and the implied total variance
$w$ (by whatever it must, to keep the price fixed).  To first order,
\[
0=\frac{\dd\big[F\,\Black(k,w)\big]}{F}
=\Black\,\dd\delta
+\partial_k\Black\,\dd k
+\partial_w\Black\,\dd w
=\big(\Black-\partial_k\Black\big)\,\dd\delta
+\partial_w\Black\,\dd w ,
\]
so
\[
\dd w
=-\frac{\Black-\partial_k\Black}{\partial_w\Black}\,\dd\delta .
\]
The numerator collapses --- this is the pleasant step:
\[
\Black-\partial_k\Black
=\PhiN(d_+)-e^k\PhiN(d_-)+e^k\PhiN(d_-)
=\PhiN(d_+),
\]
and inserting $\partial_w\Black$ --- the price's sensitivity to
total variance, the object desks call the (Black) \emph{vega} when
quoted per unit of volatility ---
\[
\dd w=-\frac{2\sqrt w\;\PhiN(d_+)}{\phin(d_+)}\,\dd\delta .
\]
Convert to volatility: $\sigma=\sqrt{w/\tau}$, so
$\dd\sigma=\dd w/(2\sqrt{w\tau})$, and the $\sqrt w$ cancels:
\begin{equation}
\frac{\dd\sigma}{\dd\delta}\bigg|_{\text{call}}
=-\frac{\PhiN(d_+)}{\phin(d_+)\,\sqrt\tau},
\qquad
\frac{\dd\sigma}{\dd\delta}\bigg|_{\text{put}}
=+\frac{\PhiN(-d_+)}{\phin(d_+)\,\sqrt\tau}.
\label{eq:fwddsig}
\end{equation}
The put case is the same computation run on the normalized put
$P=\Black+e^k-1$ (parity again): its numerator is
$P-\partial_kP=\Black-\partial_k\Black-1=\PhiN(d_+)-1=-\PhiN(-d_+)$,
which flips the sign and swaps $\PhiN(d_+)$ for $\PhiN(-d_+)$.

Before the figure, read \eqref{eq:fwddsig} at the money, where
$k=0$ and $d_+=\sqrt w/2$ is small: both $\PhiN(d_+)$ and
$\PhiN(-d_+)$ are close to $\tfrac12$, and
$\phin(d_+)\approx\phin(0)=0.399$.  A forward misread \emph{up} by
$\delta$ therefore moves the two sides in opposite directions ---
call vols down, put vols up, each by about
$\tfrac12\delta/(0.4\sqrt\tau)$.  Now recall how a smile is
assembled from quotes: at each strike the market's liquid side is the
\emph{out-of-the-money} (OTM) option --- the put below the money, the
call above --- so the drawn curve switches sides at $k=0$, and the
opposite moves become a \emph{jump} there, of size
\begin{equation}
\text{ATM gap}
\;\approx\;\frac{\delta}{\phin(0)\sqrt\tau}
\;\approx\;\frac{2.5\,\delta}{\sqrt\tau} .
\label{eq:fwdgap}
\end{equation}
The intuition for the two directions is worth one sentence each.
Thinking the forward is higher makes the call look closer to the
money than it is, and a closer-to-the-money option carrying the same
price must carry \emph{less} volatility; the same misreading pushes
the put deeper out of the money, and a deeper option at an unchanged
price must carry \emph{more}.

\Cref{fig:fwdskew} draws both halves of the argument on the running
node's fitted smile --- the frozen SPY December 2026 fit of the
earlier chapters, $\delta$ set to \MacFwdSkewBumpBp\,bp.  Panel~(a)
is the closed form \eqref{eq:fwddsig} across the quoted span, scaled
to vol basis points per \MacFwdSkewBumpBp\,bp of forward error: the
put side (positive curve) and call side (negative curve) approach the
money from opposite signs, and both fade toward the wings --- far
from the money an option's price responds weakly to the forward
relative to how strongly it responds to volatility, so the damage
concentrates exactly where the book is most liquid.  Panel~(b) is the
measurement: every out-of-the-money quote on the node re-implied at
the bumped forward, minus its true volatility.  The orange curve is
the measurement; it lands on the dotted linearization so exactly
that the two read as one curve --- and it shows the two symptoms a
desk would actually see.  At the money, the advertised jump:
\MacFwdSkewGapBp\ vol basis points of put--call gap for
\MacFwdSkewBumpBp\ basis points of forward error ---
\eqref{eq:fwdgap} with the node's numbers.  Away from the money, a tilt: the put wing lifted by
\MacFwdSkewPutTenBp\,bp at $k=-0.10$, the call wing dropped by
\MacFwdSkewCallTenBp\,bp at $k=+0.10$.  Lifted left wing, lowered
right wing --- that is precisely what a steepening of the skew looks
like.  A desk that trusts its forward would read this picture as the
market repricing crash risk; the correct reading is a data error
\$\MacFwdLineGapDollars\ wide.

\begin{figure}[htbp]
\centering
\paperfig{\textwidth}{fig_fwd_skew.pdf}
\caption{A forward error of \MacFwdSkewBumpBp\,bp on the running
node's smile.  (a)~The closed form \eqref{eq:fwddsig}: put side up,
call side down, both largest near the money.  (b)~Every OTM quote
re-implied at the bumped forward (orange; dotted: the linearization):
an ATM put--call gap of \MacFwdSkewGapBp\ vol bp plus a wing tilt
that reads exactly as a skew move.}
\label{fig:fwdskew}
\end{figure}

Now close the loop the chapter opened in \cref{sec:fwdline}.  There,
the naive whole-chain regression on the real (American) node missed
the resolved forward by \MacFwdLineGapBp\,bp --- five dollars,
invisible in price space.  Push that error through
\eqref{eq:fwdgap}: it imprints an ATM put--call gap of about
\MacFwdSkewNaiveGapBp\ vol basis points --- nearly three volatility
\emph{points} of artificial structure at the most liquid point of the
surface, dwarfing the few-basis-point fit errors Part~I worked to
achieve.  This is why the forward cannot be an afterthought: every
error in it is repaid, amplified, at the money.  And it is why the
early-exercise bow of \cref{fig:fwdline}(b) must be removed rather
than tolerated --- the task Chapter~7 takes up.

One tool from this section travels forward.  Equation
\eqref{eq:fwddsig} converts \emph{any} carry-side error into smile
units: a rate error, a dividend error, or --- next section --- a
borrow error, each reaching the smile through the forward it
misplaces.  The chapter's last question is whether that final
component of carry can be measured well enough to be worth carrying
at all.
